% from ega0-18-fr.tex
\oldpage[0\textsubscript{IV-1}]{60}

la relation (18.3.2.1). Si $\mathbf C_B$ désigne la catégorie des
$B$-bimodules, on voit qu'on a ainsi défini un \emph{foncteur covariant}
$T:\mathbf C_B\to\mathbf{Ens}$.
\end{env}

\begin{env}[18.3.3]
\label{0.18.3.3-fr}
Considérons maintenant une famille $(L_\alpha)_{\alpha\in I}$ de
$B$-bimodules et leur produit $L=\prod_{\alpha\in I}L_\alpha$~; les
projections $\operatorname{pr}_\alpha:L\to L_\alpha$ définissent des
applications
\[
  T(\operatorname{pr}_\alpha):T(L)\to T(L_\alpha),
\]
d'où une application canonique
\begin{equation}
\label{0.18.3.3.1-fr}
\tag{18.3.3.1}
  \prod_\alpha T(\operatorname{pr}_\alpha):T(L)
  \longrightarrow\prod_{\alpha\in I}T(L_\alpha).
\end{equation}
Nous allons voir que cette application est \emph{bijective}. En effet, pour
tout $\alpha\in I$, soit $E_\alpha$ une $A$-extension de $B$ par $L_\alpha$,
et soit $F$ le \emph{produit fibré} des $E_\alpha$ sur $B$ (18.1.3)~; il est
immédiat que $F$ est une $A$-extension de $B$ par
$L=\prod_{\alpha\in I}L_\alpha$ et que si l'on remplace chaque $E_\alpha$ par
une $A$-extension $A$-équivalente $E'_\alpha$, le produit fibré $F'$ des
$E'_\alpha$ sur $B$ est $A$-équivalente à $F$. On a ainsi défini une
application $\prod_{\alpha\in I}T(L_\alpha)\to T(L)$ et il est clair que cette
application est réciproque de (18.3.3.1), d'où notre assertion. On identifiera
canoniquement $T(L)$ à $\prod_{\alpha\in I}T(L_\alpha)$ par l'application
(18.3.3.1). On vérifie en outre immédiatement que si
$(L'_\alpha)_{\alpha\in I}$ est une seconde famille de $B$-bimodules et, pour
chaque $\alpha$, $w_\alpha:L_\alpha\to L'_\alpha$ un homomorphisme de
$B$-bimodules, alors, en posant
$w=\prod_\alpha w_\alpha:L\to L'=\prod_{\alpha\in I}L'_\alpha$, $T(w)$
s'identifie à $\prod_{\alpha\in I}T(w_\alpha)$ quand on fait l'identification
précédente.
\end{env}

\begin{env}[18.3.4]
\label{0.18.3.4-fr}
Cela étant, pour un $B$-bimodule $L$, l'addition est un homomorphisme
$s:L\times L\to L$ de $B$-bimodules, et il en est de même de la symétrie
$t:L\to L$ de la loi additive de $L$. On en déduit une loi de composition
\[
  T(s):T(L)\times T(L)\longrightarrow T(L)
\]
dans $T(L)$ en vertu de (18.3.3), et cette loi est une loi de \emph{groupe
commutatif} dont $T(t)$ est la symétrie, comme il résulte de la définition d'un
objet en groupes au moyen de diagrammes commutatifs
($\mathbf 0_{\mathrm{III}}$, 8.2.5 et 8.2.6). Nous désignerons par
$\operatorname{Exan}_A(B,L)$ le groupe commutatif ainsi défini et nous dirons
que c'est le \emph{groupe des classes de $A$-extensions de $B$ par $L$}.
\end{env}

\begin{env}[18.3.5]
\label{0.18.3.5-fr}
Désignons par $\mathbf K$ la catégorie dont les objets sont les triplets
$(A,B,L)$ où $A$ est un anneau, $B$ un $A$-anneau et $L$ un $B$-bimodule~;
les \emph{morphismes} de cette catégorie sont les triplets $(u,v,w)$ où
$u:A'\to A$ et $v:B'\to B$ sont deux homomorphismes d'anneaux rendant
commutatif le carré de gauche du diagramme
\[
  \xymatrix@C=2pc@R=1.5pc{
    A\ar[r] & B & L\ar[d]^w\\
    A'\ar[u]^u\ar[r] & B'\ar[u]^v & L'
  }
\]
où les flèches horizontales sont les homomorphismes structuraux~; enfin
$w:L\to L'$ est un homomorphisme de groupes commutatifs tel que l'on ait
$w(v(b')z)=b'w(z)$ et

